\documentclass[11pt,a4paper]{article}

\usepackage[T1]{fontenc}
\usepackage[utf8]{inputenc}
\usepackage[ngerman,english]{babel}
\usepackage[a4paper,margin=2.6cm]{geometry}

\usepackage{amsmath,amssymb,amsthm,mathtools}
\usepackage{bm}
\usepackage{physics}
\usepackage{graphicx}
\usepackage{booktabs}
\usepackage{enumitem}
\usepackage{hyperref}
\usepackage{microtype}
\usepackage{tikz}

\hypersetup{
	colorlinks=true,
	linkcolor=blue!50!black,
	citecolor=blue!50!black,
	urlcolor=blue!50!black
}

\title{
	A Reduced BM Four-Frame from Time, Space, Energy and Momentum\\
	\large A Spectral-Geometric Perspective on Prime Structures
}

\author{Thomas Hoffbauer}

\date{\today}

\begin{document}
	
	\maketitle
	
	\begin{abstract}
		
		We introduce a reduced four–dimensional framework of the Bamberg
		Model (BM) consisting of four distinguished directions
		$(A,C,e,b)$ interpreted as time, space, energy and momentum.
		Within this setting a Minkowski-type structure arises naturally.
		
		We show that this four-frame admits a Clifford structure and
		therefore defines a Dirac-type operator
		
		\[
		D_{\mathrm{BM}}=
		\Gamma_A A+\Gamma_C C+\Gamma_e e+\Gamma_b b .
		\]
		
		The square of this operator reproduces the invariant
		
		\[
		D_{\mathrm{BM}}^2
		=
		A^2-C^2+e^2-b^2 .
		\]
		
		Motivated by the Hilbert–Pólya idea we explore the possibility
		that the imaginary parts of the nontrivial Riemann zeros may
		appear as spectral values of such an operator.
		
		The work is presented as a heuristic program connecting
		arithmetical structures of primes with spectral geometry.
		
	\end{abstract}
	
	\tableofcontents
	
	\section{Introduction}
	
	The distribution of prime numbers and the spectral properties
	of the Riemann zeta function are among the deepest open problems
	of mathematics.
	
	The Hilbert–Pólya conjecture proposes that the imaginary parts
	of the nontrivial zeros
	
	\[
	\rho_n = \frac12 + i\gamma_n
	\]
	
	arise as eigenvalues of a self-adjoint operator.
	
	At the same time spectral geometry has shown that geometric
	structures can be encoded in the spectrum of Dirac operators.
	
	The present work proposes a reduced geometric framework,
	referred to as the Bamberg Model (BM), based on four directions
	
	\[
	(A,C,e,b)
	\]
	
	interpreted as time, space, energy and momentum.
	
	The goal is not to prove the Riemann hypothesis but to
	formulate a consistent spectral-geometric setting linking
	arithmetical prime structures with Dirac-type operators.
	
	\section{The BM Four-Frame}
	
	We introduce four distinguished directions
	
	\[
	(A,C,e,b).
	\]
	
	They are interpreted as
	
	\begin{center}
		\begin{tabular}{cc}
			Symbol & Interpretation \\
			\toprule
			$A$ & time direction \\
			$C$ & spatial direction \\
			$e$ & energy coordinate \\
			$b$ & momentum coordinate
		\end{tabular}
	\end{center}
	
	The first pair forms a reduced space-time structure,
	while the second pair corresponds to energy-momentum
	variables.
	
	\section{Reduced Space-Time Structure}
	
	Assuming $A$ to be time-like and $C$ space-like,
	a Minkowski-type interval arises naturally:
	
	\[
	ds^2 = c_M^2\,dA^2 - dC^2 .
	\]
	
	This defines an effective light cone
	
	\[
	|v|=\left|\frac{dC}{dA}\right| \le c_M .
	\]
	
	\section{Energy–Momentum Structure}
	
	The energy-momentum analogue of the Minkowski interval is
	
	\[
	e^2 - c_M^2 b^2.
	\]
	
	This leads to a relativistic invariant
	
	\[
	e^2 - c_M^2 b^2
	=
	m_{\mathrm{BM}}^2 c_M^4 .
	\]
	
	\section{Clifford Structure}
	
	The four BM directions naturally define a Clifford algebra
	
	\[
	Cl(1,3)
	\]
	
	with generators $\Gamma_\mu$ satisfying
	
	\[
	\{\Gamma_\mu,\Gamma_\nu\}=2g_{\mu\nu}.
	\]
	
	\section{BM Dirac Operator}
	
	We define a Dirac-type operator
	
	\[
	D_{\mathrm{BM}}
	=
	\Gamma_A A
	+
	\Gamma_C C
	+
	\Gamma_e e
	+
	\Gamma_b b .
	\]
	
	Using the Clifford relations one obtains
	
	\[
	D_{\mathrm{BM}}^2
	=
	A^2 - C^2 + e^2 - b^2 .
	\]
	
	In the light-like case
	
	\[
	A^2-C^2=0
	\]
	
	this reduces to
	
	\[
	D_{\mathrm{BM}}^2
	=
	e^2-b^2 .
	\]
	
	\section{Prime Structure}
	
	The multiplicative group
	
	\[
	(\mathbb Z/12\mathbb Z)^\times
	=
	\{1,5,7,11\}
	\]
	
	forms the Klein four group
	
	\[
	V_4 .
	\]
	
	Prime quadruples of these residue classes appear naturally
	in the BM framework.
	
	\section{Numerical Experiments}
	
	To explore the BM framework numerically we consider
	prime quadruples
	
	\[
	(p_1,p_2,p_3,p_4)
	\]
	
	with residues
	
	\[
	\{1,5,7,11\} \pmod{12}.
	\]
	
	Coordinates are defined as
	
	\[
	C = \log(p_1p_2p_3p_4)
	\]
	
	and
	
	\[
	b = \log(p_1^2+p_2^2+p_3^2+p_4^2).
	\]
	
	Combined with Riemann zero ordinates $\gamma_n$
	we obtain BM coordinates
	
	\[
	(A,C,e,b).
	\]
	
	\section{BM Dirac Spectrum}
	
	\begin{figure}[h]
		\centering
		\begin{tikzpicture}
			
			\draw[->] (0,0) -- (8,0);
			\draw[->] (0,0) -- (0,5);
			
			\foreach \x in {0.5,1,...,7.5}
			\draw[blue] (\x,{0.5*\x+0.3*rand}) circle (0.05);
			
			\node at (4,-0.4) {$\gamma$};
			\node at (-0.3,4) {$|\lambda|$};
			
		\end{tikzpicture}
		
		\caption{Illustration of Dirac-type spectral scaling
			$\lambda\sim\gamma$.}
		
	\end{figure}
	
	\section{Prime Shell Structure}
	
	\begin{figure}[h]
		\centering
		
		\begin{tikzpicture}
			
			\draw (0,0) circle (1);
			\draw (0,0) circle (2);
			\draw (0,0) circle (3);
			
			\node at (1,0) {$1$};
			\node at (0,2) {$5$};
			\node at (-2,0) {$7$};
			\node at (0,-3) {$11$};
			
		\end{tikzpicture}
		
		\caption{Schematic illustration of V4 prime shells
			in the BM framework.}
		
	\end{figure}
	
	\section{Discussion}
	
	The BM construction should be understood as a heuristic
	framework.
	
	The numerical experiments do not establish a direct spectral
	identity between the BM Dirac operator and the Riemann zeros.
	
	However the emergence of Dirac-like spectral structures
	suggests that the four-frame $(A,C,e,b)$ may provide a useful
	geometric interpretation of arithmetic patterns.
	
	\section{Outlook}
	
	Future work may include
	
	\begin{itemize}
		\item construction of a self-adjoint BM Dirac operator
		\item deeper analysis of the spectral distribution
		\item connections to noncommutative spectral geometry
	\end{itemize}
	
	\begin{thebibliography}{9}
		
		\bibitem{Edwards}
		H. M. Edwards,
		\textit{Riemann's Zeta Function},
		Dover.
		
		\bibitem{BerryKeating}
		M. Berry, J. Keating,
		The Riemann zeros and eigenvalue asymptotics.
		
		\bibitem{Connes}
		A. Connes,
		Trace formula in noncommutative geometry.
		
	\end{thebibliography}
	
\end{document}